%
% drehung.tex -- Drehungen
%
% (c) 2020 Prof Dr Andreas Müller, Hochschule Rapperswil
%
% !TEX root = ../../buch.tex
% !TEX encoding = UTF-8
%

\section{Drehungen\label{qa:section:drehungen}}
\kopfrechts{Drehungen}
\subsubsection{Eulersche Formel}
Bevor Drehungen durchgeführt werden können, müssen die vier 
Dimensionen einem Zweck zugeordnet werden. Das triviale  Quaternion 
ist $1 + 0i + 0j + 0k$. Eine Multiplikation mit diesem Quaternion 
verändert keine andere Zahl. Die drei imaginären Achsen $i,j,k$ 
werden den drei Raumachsen $x,y,z$ zugeordnet. Damit kann eine 
Drehung um die $x$-Achse mit dem Quaternion $0 + 1i + 0j + 0k$ 
durchgeführt werden. Analog dazu kann eine Drehung um die $y$-Achse 
mit dem Quaternion $0 + 0i + 1j + 0k$ und eine Drehung um die 
$z$-Achse mit dem Quaternion $0 + 0i + 0j + 1k$ durchgeführt werden.

Der Realteil (und ein Faktor im Imaginärteil) geben die Stärke der 
Drehung an. Dafür wird die Euler'sche Drehformel in der Form
\begin{equation*}
    e^{i\theta} = \cos(\theta) + i\sin(\theta)
\end{equation*}
verwendet, welche bei den komplexen Zahlen eine Drehung um $\theta$ 
beim Nullpunkt in der komplexen Ebene beschreibt. Für Quaternionen 
wird die Formel erweitert, sodass eine Drehung um eine beliebige 
Achse durchgeführt werden kann. Die Achse wird durch einen 
Einheitsvektor beschrieben, welcher die Richtung der Achse angibt. 
Der Winkel $\theta$ gibt, identisch wie bei den komplexen Zahlen, 
die Stärke der Drehung an. Damit kann ein Quaternion für eine 
Drehung um eine beliebige Achse mit der Formel
\begin{equation*}
    q = e^{\theta(xi + yj + zk)} 
    = \cos \left(\theta\right) 
    + \frac{xi + yj + zk}{|x + y + z|}\sin\left(\theta\right)
\end{equation*}
berechnet werden. Der Bruch vor dem Sinus-Term stellt sicher, dass 
der Imaginärteil des Quaternions ein Einheitsvektor ist. Die Formel 
ist jedoch in dieser Form noch nicht anwendbar, wie im nächsten 
Beispiel gezeigt wird.

\subsubsection{Beispiel: Drehung um die $z$-Achse}
Sei ein Punkt $P_1$ auf der $z$-Achse mit den Koordinaten $(0,0,1)
$. Mit Quaternionen lässt sich $P_1$ wie folgt darstellen: $P_1 = 0
+ 0i + 0j + 1k$. Dieser Punkt soll um die $z$-Achse um 90° (oder 
$\frac{\pi}{2}$) gedreht werden. Das erwartete Ergebnis ist 
wiederum $(0,0,1)$, da eine Drehung um die eigene Achse keine 
Wirkung hat. Das Quaternion für diese Drehung wird mit 
\begin{equation*}
    q = e^{\frac{\pi}{2}(0i + 0j + 1k)} 
    = \cos\left(\frac{\pi}{2}\right) 
    + \frac{0i + 0j + 1k}{|0 + 0 + 1|}
    \sin\left(\frac{\pi}{2}\right) 
    = k
\end{equation*}
berechnet. Nun soll die Drehung durchgeführt werden. Dafür wird das 
Quaternion $q$ mit dem Punkt $P_1$ von links multipliziert:
\begin{equation*}
    P_{1L} = q \cdot P_1 = k \cdot k = k^2 = -1
\end{equation*}
Das Ergebnis ist $-1$, was dem Punkt $(-1,0,0,0)$ entspricht. Statt 
dass der Punkt auf der $z$-Achse geblieben ist, wurde er auf den 
Ursprung verschoben mit Länge 1 in der vierten Dimension, welches 
nicht dem erwarteten Ergebnis entspricht. 

Vierdimensionale Zahlen lassen sich nicht leicht visualisieren. 
Eine Möglichkeit ist, die imaginären Achsen in einem 
dreidimensionalen Koordinatensystem und die vierte Dimension 
zusammen mit der $z$-Achse in einem zweidimensionalen 
Koordinatensystem darzustellen. Die Animation zu dieser Drehung 
kann über den QR Code \ref{qa:fig:QRDrehung} abgespielt werden.
\begin{figure}
    \centering
    \qrcode[height=3cm]{https://youtu.be/7kv7XFGp4ns}
    \caption{QR-Code zum Video mit der Animation der Drehung um die
    $z$-Achse.}
    \label{qa:fig:QRDrehung}
\end{figure}
Der rote, grüne und blaue Pfeil in der Animationzeigen jeweils in 
die $x,y,z$-Achse und repräsentieren die imaginären $i,j,k$-Achsen. 
Der gelbe Pfeil zeigt in die reele Achse und ist orthogonal zu 
allen anderen Achsen. In anderen Worten: Man muss sich vorstellen, 
dass der gelbe Pfeil rechtwinklig steht zu allen anderen Pfeilen, 
in der vierten Dimension. Der Punkt $P_1$ liegt nun tatsächlich im 
Ursprung, mit -1 in der reelen Achse.

Um den Punkt tatsächlich um 90° um die $z$-Achse im 
Gegenuhrzeigersinn zu drehen, muss die Rolle der Reihenfolge 
analysiert werden. Dafür wird ein neuer Beispielpunkt $P_2 = (0,1,1)
$ verwendet, welcher nach einer Drehung auch den Ort wechselt. Eine 
Multiplikation von Links würde
\begin{equation*}
    P_{2L} = q P_2 
    = k \cdot (0 + 0i + 1j + 1k) 
    = k (j + k) 
    = kj + k^2 
    = -i - 1
\end{equation*}
ergeben, welches dem Punkt $(-1,0,0)$ im 3D-Raum entspricht. Der 
Punkt wurde also um 90° um die $z$-Achse im Gegenuhrzeigersinn 
gedreht, jedoch wurde der $z$-Anteil des Punktes auf 0 reduziert. 
Der Realteil beträgt -1.

Von Rechts würde die Multiplikation
\begin{equation*}
    P_{2R} 
    = P_2 q 
    = (0 + 0i + 1j + 1k) k 
    = jk + k^2 
    = i - 1
\end{equation*}
lauten, welches dem Punkt $(1,0,0)$ entspricht. Von Rechts wurde 
der Punkt also um 90° in umgekehrter Drehrichtung um die $z$-Achse 
gedreht. Auch hier wurde der $z$-Anteil des Punktes auf 0 
reduziert, mit Realteil -1.

Die Multiplikation von Links und Rechts betrifft also primär die 
Richtung der Drehung. Es ist zu beachten, dass für beide 
Multiplikationen das selbe Phänomen auftritt, dass der $z$-Anteil 
des Punktes auf 0 reduziert und der Realteil auf -1 gesetzt wird. 
Anderst ausgedrückt: Die entstandene Verzerrung ist sowohl bei 
Links als auch bei Rechts identisch. Im nächsten Abschnitt wird 
gezeigt, wie dieser Nebeneffekt elegant gelöst werden kann.
\subsubsection{Die Drehformel}%
Für die effektive Drehung werden die folgenden Eigenschaften der 
Quaternionen verwendet:
\begin{itemize}
    \item Multiplikation von Links: Gegenuhrzeigersinn mit 
    Verzerrung
    \item Multiplikation von Rechts: Uhrzeigersinn mit identischer 
    Verzerrung
    \item Inverse eines Quaternions: $q^{-1}$ macht die Drehung 
    rückgängig, sodass $q q^{-1} = 1$
    \item Bei Einheitsquaternionen gilt: $q^{-1} = q*$, wobei $q*$ 
    das konjugierte Quaternion ist
\end{itemize}
Für die verzerrungsfreie Drehung wird zuerst der Punkt $P$ mit dem 
Quaternion $q$ von Links multipliziert, was in unserem Beispiel 
eine Drehung um 90° im Gegenuhrzeigersinn ergibt. Danach wird das 
Ergebnis von Rechts mit der Inverse desselben Quaternions $q^{-1}$ 
multipliziert, welches erneut eine Drehung um 90° im 
Gegenuhrzeigersinn bewirkt. Die Inverse hat dabei zwei 
entscheidende Wirkungen:
\begin{itemize}
    \item Die zweite Drehung macht die Verzerrung rückgängig
    \item Die Drehrichtung von Rechts wird umgekehrt, sodass zwei 
    aufeinanderfolgende Drehungen in gleicher Richtung erfolgen
\end{itemize}
Die Drehformel lautet also:
\begin{equation*}
    P' = q P q^{-1}
\end{equation*}
Und für Einheitsquaternionen kann die Formel auch in der Form
\begin{equation*}
    P' = q P q*
\end{equation*}
geschrieben werden. Die Berechnung der Konjugation ist deutlich einfacher als die Berechnung der Inversen, da nur die Vorzeichen der Imaginärteile geändert werden müssen.
Mit dem Beispielpunkt $P_2$ und dem Quaternion $q$ für eine Drehung um 90° um die $z$-Achse ergibt sich:
\begin{align*}
    P' &= q P_2 q^{-1} \\
    &= k (0 + 0i + 1j + 1k) k^{-1} \\
    &= k (j + k) (-k) \\
    &= kjk + kk(-k) \\
    &= -i - k^3 \\
    &= -i + k \\
    &= (-1, 0, 1) \quad \text{(ohne Realteil!)}
\end{align*}
Es ist zu beachten, dass nun eine 180° Drehung durchgeführt wurde, 
da die Multiplikation von Links und Rechts jeweils eine 90° Drehung 
bewirkt. Das heisst, um eine Drehung um 90° zu erreichen, muss der 
Winkel $\theta$ in der Berechnung des Quaternions $q$ auf $\frac
{\pi}{4}$ gesetzt werden. Um das Drehquaternion zu berechnen, wird 
die Formel
\begin{equation*}
    q = \cos\left(\frac{\theta}{2}\right) 
    + \frac{xi + yj + zk}{|x + y + z|}
    \sin\left(\frac{\theta}{2}\right)
\end{equation*}
so umgestellt, dass der Winkel $\theta$ halbiert wird.